\section{Архітектура та компоненти системи}

Для проведення експериментального оцінювання якості перекладу мовлення
було обрано каскадну архітектуру (див. підрозділ~\ref{subsec:cascade}), у якій етапи
розпізнавання мовлення та машинного перекладу працюють послідовно:
вхідний аудіосигнал українською $x(t)$ перетворюється на текст $\hat{y}$
(ASR), який потім перекладається на англійську $\hat{z}$ (MT).

Компонент синтезу мовлення (TTS) не включено до pipeline, оскільки
дослідження зосереджене на оцінюванні якості саме текстового перекладу.
Аналіз впливу TTS-компоненту на загальну якість може бути предметом
подальших досліджень повного pipeline перекладу мовлення.

\\
Каскадний підхід обрано з огляду на такі переваги:
\begin{itemize}[nosep]
    \item \textit{модульність} --- можливість замінювати окремі компоненти
    та порівнювати різні MT-системи при одному й тому ж ASR-виході;
    \item \textit{інтерпретованість} --- якість кожного етапу оцінюється
    окремо (WER/CER для ASR, BLEU/BERTScore/METEOR для MT), що дозволяє
    локалізувати джерела помилок;
    \item \textit{аналіз каскадного поширення помилок} --- можливість
    дослідити, яким чином помилки розпізнавання мовлення впливають на
    якість фінального перекладу;
    \item \textit{відтворюваність} --- використання відкритих моделей
    гарантує повну відтворюваність експерименту.
\end{itemize}

Альтернативою каскадному підходу є end-to-end моделі, зокрема
\newline SeamlessM4T~\cite{seamlessm4t}, що здійснюють переклад мовлення напряму,
без проміжного текстового подання. Однак якість таких моделей для мовної
пари українська--англійська на момент проведення дослідження поступається
каскаду з окремих компонентів~\cite{seamlessm4t}. Крім того, end-to-end
підхід унеможливлює покомпонентний аналіз помилок, що є одним із завдань
цієї роботи.

Модель Whisper також має вбудований режим прямого перекладу
\newline(\texttt{task="translate"}), який виконує speech-to-English-text без
проміжного етапу. Проте якість такого перекладу для пари uk$\to$en значно
нижча за каскад Whisper ASR + спеціалізована MT-модель, оскільки обсяг
тренувальних даних для прямого перекладу з української є значно меншим.

\subsection{ASR-компонент: Whisper medium}

Для автоматичного розпізнавання мовлення обрано модель OpenAI Whisper
medium~\cite{whisper} --- мультимовну ASR-модель, натреновану на
приблизно 680\,000 годин слабко-розміченого аудіо з мережі. Модель
підтримує 99 мов, включно з українською, та застосовує zero-shot підхід,
що не потребує дотренування для конкретної мови чи домену.

Whisper medium містить 769~млн параметрів та має оптимальний
баланс між якістю розпізнавання та обчислювальними вимогами. За даними
Whisper Leaderboard на платформі Hugging Face, WER моделі medium для
українського мовлення становить приблизно 10--15\% на чистому мовленні.

Було розглянуто такі альтернативи:
\begin{itemize}[nosep]
    \item \textit{Whisper large-v3} --- має вищу точність (WER на
    2--3\% нижчий), проте потребує значно більших обчислювальних ресурсів;
    \item \textit{Whisper turbo} --- працює швидше,
    проте має нижчу точність на малоресурсних мовах;
    \item \textit{wav2vec~2.0 (uk)} --- потребує fine-tuning на цільових
    даних та має менший обсяг доступних тренувальних даних для української;
    \item \textit{Google Speech-to-Text} --- комерційна система з
    обмеженою відтворюваністю результатів.
\end{itemize}

Whisper medium обрано як оптимальний варіант для дослідження з урахуванням
наявних обчислювальних ресурсів. Для практичної
ефективності використано реалізацію Faster-Whisper на базі CTranslate2,
що дає приблизно 4-кратне прискорення інференсу~\cite{ctranslate2}
при збереженні якості розпізнавання завдяки INT8-квантизації.

\subsection{MT-компонент: OPUS-MT}

Для етапу машинного перекладу обрано спеціалізовану NMT-модель
Helsinki-NLP/opus-mt-uk-en~\cite{opus_mt}, побудовану на архітектурі
Transformer у фреймворку MarianNMT~\cite{marian_nmt}. Модель натренована
на корпусі OPUS, що включає CCMatrix, WikiMatrix, Europarl та інші
паралельні дані. Розмір моделі --- приблизно 74~млн параметрів, максимальна
довжина контексту --- 512 токенів. Очікуваний BLEU на бенчмарку
FLORES-101~\cite{flores} для uk$\to$en становить 25--30 залежно від домену.
Модель повністю відкрита (ліцензія MIT), працює локально та має
високу швидкість інференсу.

Вибір OPUS-MT як основного MT-компонента зумовлений кількома факторами:
\begin{itemize}[nosep]
    \item \textit{відкритість} --- повна відтворюваність результатів без
    залежності від комерційних API;
    \item \textit{спеціалізація} --- модель натренована саме для пари
    uk$\to$en, на відміну від універсальних систем;
    \item \textit{ефективність} --- невеликий розмір моделі дає змогу
    працювати на обмежених обчислювальних ресурсах навіть без використання GPU.
\end{itemize}

Серед альтернативних NMT-систем було розглянуто NLLB-200~\cite{nllb}
(1,3~млрд параметрів, BLEU 35--38 на FLORES для uk$\to$en), проте OPUS-MT
обрано як більш легковісний та відтворюваний baseline, що дає змогу чітко
виділити каскадний ефект без впливу масштабу моделі.

*Як додаткове дослідження аналізу якості перекладу також залучено три великі мовні
моделі (GPT-4o, Claude Sonnet 4, Gemini 2.0 Flash), опис та результати
яких наведено у підрозділі~\ref{subsec:llm_comparison}.

\subsection{Параметри експерименту}

Для кожного аудіосегменту корпусу фіксуються такі дані:
\begin{enumerate}[nosep]
    \item час виконання ASR (у секундах);
    \item текст ASR-виходу (розпізнана транскрипція українською);
    \item час виконання MT (у секундах);
    \item текст MT-виходу (переклад англійською);
    \item загальна затримка pipeline (сума часу ASR та MT).
\end{enumerate}

Параметри ASR-компоненту наведено у таблиці~\ref{tab:whisper} (підрозділ~\ref{subsec:corpus_prep}).
Для OPUS-MT використано параметри за замовчуванням бібліотеки Hugging Face
Transformers.

\subsection{Параметри експериментального прогону}

Для обробки корпусу було розроблено автоматизований \textit{pipeline}, що
поєднує два ключові етапи: автоматичне розпізнавання мовлення (ASR) та
машинний переклад (MT). Обробка виконувалась на персональному комп'ютері
з процесором та без виділеного GPU, що дає змогу оцінити продуктивність
системи в умовах обмежених обчислювальних ресурсів.

\subsubsection{Використані моделі}

Для ASR використано модель Whisper medium (параметри наведено в
таблиці~\ref{tab:whisper}). Для MT --- модель
Helsinki-NLP/opus-mt-uk-en (характеристики наведено в
підрозділі~4.2). Обидві моделі детально описано
в підрозділах~4.1--4.2.

\subsubsection{Результати прогону}

Характеристики експериментального корпусу детально описано в розділі~\ref{sec:corpus}
(таблиця~\ref{tab:corpus}).

Весь корпус було оброблено в автоматичному режимі із збереженням:
\begin{itemize}
    \item тексту транскрипції (ASR) українською мовою;
    \item тексту машинного перекладу (MT) англійською мовою;
    \item часових міток (timestamps) для кожного сегмента мовлення;
    \item метрик латентності для ASR та MT етапів.
\end{itemize}

\begin{table}[h]
\centering
\caption{Метрики продуктивності pipeline}
\label{tab:latency}
\begin{tabular}{|l|r|r|}
\hline
\textbf{Етап} & \textbf{Час (с)} & \textbf{RTF} \\
\hline
ASR (Whisper medium) & 361,3 & 0,18$\times$ \\
MT (OPUS-MT) & 156,7 & --- \\
\hline
\textbf{Загалом} & \textbf{518,0} & \textbf{0,26$\times$} \\
\hline
\end{tabular}
\end{table}

Коефіцієнт RTF (\textit{Real-Time Factor}) показує відношення часу обробки
до тривалості аудіо. Значення RTF = 0,18$\times$ для ASR означає, що система
обробляє аудіо приблизно в 5,5 разів швидше за реальний час навіть без GPU.
Загальний RTF = 0,26$\times$ (ASR + MT) підтверджує придатність системи
для офлайн-обробки великих корпусів.

Результати збережено у структурованому форматі JSON, що містить для кожного
аудіофайлу:
\begin{itemize}
    \item масив сегментів із часовими мітками \texttt{start}/\texttt{end};
    \item оригінальний текст (\texttt{text\_uk}) та переклад (\texttt{text\_en});
    \item метрики латентності (\texttt{asr\_time\_sec}, \texttt{mt\_time\_sec}).
\end{itemize}

% Підключення підрозділів
\subsection{Система метрик оцінювання якості перекладу}
\label{subsec:metrics}

Якість перекладу оцінено автоматичними метриками та ручною експертною
оцінкою. Автоматичні метрики вимірюють формальну відповідність референсному
тексту, тоді як експертна оцінка враховує семантичну адекватність та
природність мовлення.

\subsubsection{Автоматичні метрики}

\paragraph{BLEU (Bilingual Evaluation Understudy).}
Метрика BLEU~\cite{papineni_bleu} є однією з найбільш поширених метрик для
оцінки машинного перекладу. Вона базується на порівнянні $n$-грам
(послідовностей із $n$ слів) між машинним перекладом та одним або кількома
референсними перекладами.

Формула обчислення BLEU:
\begin{equation}
\text{BLEU} = BP \cdot \exp\left(\sum_{n=1}^{N} w_n \log p_n\right)
\end{equation}
де $p_n$ — модифікована точність (precision) для $n$-грам, $w_n$ — вагові
коефіцієнти (зазвичай $w_n = 1/N$), $N$ — максимальний порядок $n$-грам
(стандартно $N = 4$), а $BP$ — штраф за стислість (brevity penalty):
\begin{equation}
BP = \begin{cases}
1 & \text{якщо } c > r \\
e^{(1 - r/c)} & \text{якщо } c \leq r
\end{cases}
\end{equation}
де $c$ — довжина машинного перекладу, $r$ — ефективна референсна довжина.

Модифікована точність для $n$-грам обчислюється як:
\begin{equation}
p_n = \frac{\sum_{\text{C} \in \{\text{Candidates}\}} \sum_{\text{n-gram} \in C} Count_{clip}(\text{n-gram})}{\sum_{\text{C} \in \{\text{Candidates}\}} \sum_{\text{n-gram} \in C} Count(\text{n-gram})}
\end{equation}

BLEU обрано як стандарт де-факто для порівняння систем машинного
перекладу~\cite{post_bleu}: метрика дає швидке та відтворюване оцінювання
і корелює з людською оцінкою на великих корпусах.

\paragraph{BERTScore.}
На відміну від BLEU, метрика BERTScore~\cite{bert_score} оцінює семантичну
подібність на рівні контекстуалізованих векторних представлень слів,
отриманих за допомогою попередньо навченої мовної моделі BERT.

Для кожного токена $x_i$ в машинному перекладі та токена $\hat{x}_j$ в
референсі обчислюється косинусна подібність їхніх контекстуалізованих
ембедингів:
\begin{equation}
\text{sim}(x_i, \hat{x}_j) = \frac{\mathbf{x}_i^\top \hat{\mathbf{x}}_j}{\|\mathbf{x}_i\| \|\hat{\mathbf{x}}_j\|}
\end{equation}

На основі цих подібностей обчислюються три компоненти:
\begin{align}
P_{BERT} &= \frac{1}{|x|} \sum_{x_i \in x} \max_{\hat{x}_j \in \hat{x}} \text{sim}(x_i, \hat{x}_j) \\
R_{BERT} &= \frac{1}{|\hat{x}|} \sum_{\hat{x}_j \in \hat{x}} \max_{x_i \in x} \text{sim}(x_i, \hat{x}_j) \\
F_{BERT} &= 2 \cdot \frac{P_{BERT} \cdot R_{BERT}}{P_{BERT} + R_{BERT}}
\end{align}
де $P_{BERT}$ — точність, $R_{BERT}$ — повнота, $F_{BERT}$ — F1-міра.

BERTScore обрано, оскільки метрика краще враховує семантичну еквівалентність
перефразувань та синонімів і має вищу кореляцію з людською оцінкою
порівняно з BLEU для оцінки окремих речень, що важливо для
сегментованого мовлення.

\paragraph{METEOR (Metric for Evaluation of Translation with Explicit ORdering).}
Метрика METEOR~\cite{banerjee_meteor} розроблена для усунення деяких недоліків
BLEU. Вона враховує не лише точний збіг слів, але й стемінг (зведення до
основи слова), синоніми (на основі WordNet) та парафрази.

METEOR обчислює гармонійне середнє між точністю та повнотою уніграм з
урахуванням штрафу за фрагментацію:
\begin{equation}
\text{METEOR} = F_{mean} \cdot (1 - Penalty)
\end{equation}
де
\begin{equation}
F_{mean} = \frac{10 \cdot P \cdot R}{R + 9P}
\end{equation}
\begin{equation}
Penalty = 0.5 \cdot \left(\frac{\#chunks}{\#unigrams\_matched}\right)^3
\end{equation}

Тут $P$ — точність, $R$ — повнота, $\#chunks$ — кількість неперервних
фрагментів зіставлених уніграм.

METEOR має кращу кореляцію з людською оцінкою на рівні речень порівняно
з BLEU, особливо для мов із багатою морфологією. Врахування синонімів та
стемінгу є важливим для оцінки перекладу публіцистичних текстів.

\subsubsection{Людська експертна оцінка}

Автоматичні метрики, попри свою ефективність, не завжди адекватно відображають
якість перекладу з точки зору людини-читача. Тому в методології передбачено
можливість експертної оцінки за двома критеріями. На практиці основну увагу
зосереджено на анотації помилок та оцінці когерентності (розділ~\ref{sec:results}), оскільки
ці методи надають детальнішу діагностичну інформацію, ніж загальні шкали
адекватності та плавності.

\paragraph{Адекватність та плавність.}
Для оцінки якості перекладу можуть застосовуватися шкали
адекватності (ступінь збереження змісту оригіналу) та плавності
(природність мовлення цільовою мовою) за 5-бальною шкалою.
У роботі основну увагу зосереджено на більш діагностичних
методах --- анотації помилок за типами та оцінці міжфразової
когерентності, що надають детальнішу інформацію для аналізу каскадного
ефекту.

\subsubsection{Класифікація помилок перекладу}

Для детального аналізу помилок застосовано базову класифікацію, описану
в підрозділі~\ref{subsec:error_class} (пропуски, спотворення, додавання), розширену додатковими
категоріями на основі фреймворку MQM~\cite{mqm} (Multidimensional Quality
Metrics):
\begin{itemize}[nosep]
    \item \textit{Помилки в іменах власних} --- неправильна транслітерація
    імен, географічних назв, абревіатур;
    \item \textit{Граматичні помилки} --- порушення граматичних норм
    цільової мови;
    \item \textit{Стилістичні помилки} --- неприродні конструкції,
    калькування з мови оригіналу.
\end{itemize}

Кожна помилка додатково класифікується за ступенем серйозності:
\begin{itemize}[nosep]
    \item \textbf{Critical} (критична) --- помилка, що спотворює основний зміст
    повідомлення або містить фактичну неточність;
    \item \textbf{Major} (суттєва) --- помилка, що ускладнює розуміння фрагменту;
    \item \textbf{Minor} (незначна) --- помилка, що не впливає на загальне
    розуміння тексту.
\end{itemize}

\subsection{Дизайн експерименту}

У даному підрозділі описано методологію проведення експериментального
дослідження якості машинного перекладу, включаючи процедуру оцінювання,
статистичні методи аналізу та забезпечення валідності результатів.

\subsubsection{Загальна схема експерименту}

Експеримент побудовано за схемою порівняння машинного перекладу з
референсним людським перекладом.
Основні етапи:

\begin{enumerate}
    \item \textbf{Підготовка корпусу:} сегментація аудіозаписів на фрагменти
    оптимальної тривалості (3--5 хвилин);
    \item \textbf{Автоматична обробка:} послідовне застосування ASR (Whisper)
    та MT (OPUS-MT) до кожного сегмента;
    \item \textbf{Створення референсу:} ручна транскрипція та підготовка
    референсних перекладів тих самих сегментів (деталі в розділі~\ref{sec:corpus});
    \item \textbf{Автоматичне оцінювання:} обчислення метрик BLEU, BERTScore,
    METEOR для кожного сегмента;
    \item \textbf{Експертне оцінювання:} оцінка адекватності та плавності,
    анотація помилок;
    \item \textbf{Статистичний аналіз:} агрегація результатів, перевірка
    гіпотез, аналіз кореляцій.
\end{enumerate}

\subsubsection{Процедура автоматичного оцінювання}

Автоматичні метрики обчислюються з використанням стандартизованих
інструментів для забезпечення відтворюваності результатів:

\begin{itemize}
    \item \textbf{BLEU:} бібліотека SacreBLEU~\cite{post_bleu} з токенізацією
    ``13a'' (стандарт WMT);
    \item \textbf{BERTScore:} модель \texttt{microsoft/deberta-xlarge-mnli}
    для англійської мови з налаштуванням \texttt{rescale\_with\_baseline=True};
    \item \textbf{METEOR:} реалізація з пакету NLTK з використанням WordNet
    для синонімів.
\end{itemize}

Оцінювання проводиться на рівні окремих сегментів, що дозволяє:
\begin{itemize}
    \item аналізувати варіативність якості перекладу;
    \item виявляти проблемні ділянки тексту;
    \item досліджувати вплив характеристик вхідного аудіо на якість.
\end{itemize}

\subsubsection{Процедура експертного оцінювання}

Експертне оцінювання зосереджено на двох ключових аспектах: анотації
помилок за типами та оцінці міжфразової когерентності. Оцінювач має
доступ до оригінальної транскрипції українською, машинного перекладу
англійською та референсного перекладу.

\paragraph{Анотація помилок.}
Для детального аналізу проводиться анотація помилок за класифікацією,
наведеною в підрозділі~\ref{subsec:error_class}. Для кожної помилки фіксується:
\begin{itemize}
    \item тип помилки (пропуск, додавання, спотворення тощо);
    \item ступінь серйозності (critical, major, minor);
    \item контекст (фрагмент тексту з помилкою).
\end{itemize}

\subsubsection{Статистичні методи аналізу}

Для аналізу отриманих результатів застосовано наступні статистичні методи:

\paragraph{Описова статистика.}
Для кожної метрики обчислюються:
\begin{itemize}
    \item середнє значення ($\bar{x}$);
    \item стандартне відхилення ($\sigma$);
    \item медіана (Me);
    \item мінімальне та максимальне значення.
\end{itemize}

\paragraph{Кореляційний аналіз.}
Для дослідження зв'язку між автоматичними метриками та людською оцінкою
обчислюється коефіцієнт кореляції Спірмена~\cite{spearman} ($\rho$),
який є більш робастним для ординальних даних:
\begin{equation}
\rho = 1 - \frac{6 \sum d_i^2}{n(n^2 - 1)}
\end{equation}
де $d_i$ — різниця між рангами $i$-го спостереження, $n$ — кількість
спостережень.

\paragraph{Перевірка статистичної значущості.}
Для перевірки значущості відмінностей між групами застосовується критерій
Манна-Уітні (U-тест) як непараметрична альтернатива t-тесту, що не вимагає
припущення про нормальний розподіл:
\begin{equation}
U = n_1 n_2 + \frac{n_1(n_1 + 1)}{2} - R_1
\end{equation}
де $n_1$, $n_2$ — розміри вибірок, $R_1$ — сума рангів першої вибірки.

Рівень значущості встановлено на $\alpha = 0.05$.

\subsubsection{Аналіз впливу помилок ASR на якість перекладу}

Окремим напрямком аналізу є дослідження каскадного ефекту — впливу помилок
автоматичного розпізнавання мовлення (ASR) на якість машинного перекладу.
Для цього:

\begin{enumerate}
    \item Обчислюється WER (Word Error Rate) між ASR-транскрипцією та
    референсною транскрипцією за формулою, наведеною в підрозділі~\ref{subsec:asr}.

    \item Проводиться кореляційний аналіз між WER/CER та метриками якості
    перекладу (BLEU, BERTScore, METEOR).

    \item Виділяються типові помилки ASR, що призводять до критичних
    помилок перекладу (зокрема, спотворення імен власних та абревіатур).
\end{enumerate}

\subsubsection{Забезпечення валідності дослідження}

Для забезпечення валідності та надійності результатів дослідження вжито
наступних заходів:

\begin{itemize}
    \item \textbf{Внутрішня валідність:} використання стандартизованих
    інструментів обчислення метрик, чітких критеріїв експертної оцінки,
    контроль за процедурою оцінювання.

    \item \textbf{Зовнішня валідність:} корпус містить реальні публічні
    промови різної тематики, що підвищує екологічну валідність результатів.

    \item \textbf{Відтворюваність:} всі параметри моделей, версії бібліотек
    та налаштування задокументовані, вихідні дані збережені у структурованому
    форматі.
\end{itemize}

\begin{table}[h]
\centering
\caption{Зведена таблиця методів оцінювання}
\label{tab:evaluation_methods}
\begin{tabular}{|l|l|l|}
\hline
\textbf{Аспект оцінки} & \textbf{Метод} & \textbf{Рівень} \\
\hline
Лексичний збіг & BLEU & Сегмент/корпус \\
Семантична подібність & BERTScore & Сегмент \\
Морфологічний збіг & METEOR & Сегмент \\
Типи помилок & Анотація помилок (MQM) & Сегмент \\
Когерентність & Експертна оцінка & Пара сегментів \\
Якість ASR & WER, CER & Сегмент \\
\hline
\end{tabular}
\end{table}

\subsection*{Висновки до розділу 4}

У четвертому розділі роботи описано методологію експериментального дослідження
автоматичного перекладу усного мовлення з української на англійську мову.
Основні результати розділу:

\begin{enumerate}

\item \textbf{Сформовано експериментальний корпус} із аудіозаписів
публічних звернень українською мовою загальною тривалістю близько 45 хвилин,
сегментованих на 12 фрагментів для оптимальної обробки моделями.
Корпус характеризується високою тематичною когерентністю та містить
типові для публіцистичного стилю мовні конструкції, імена власні та
абревіатури.

\item \textbf{Реалізовано автоматизований pipeline} обробки мовлення,
що поєднує модель автоматичного розпізнавання мовлення Whisper medium
та спеціалізовану модель машинного перекладу OPUS-MT.
Продуктивність системи (RTF = 0,26$\times$) підтверджує
її придатність для офлайн-обробки великих корпусів навіть на
обмежених обчислювальних ресурсах без використання GPU.

\item \textbf{Обґрунтовано вибір системи метрик} для оцінки якості
перекладу. Комбінація автоматичних метрик BLEU (лексичний збіг),
BERTScore (семантична подібність) та METEOR (морфологічний збіг із
урахуванням синонімів) дає повну кількісну оцінку.
Доповнення ручною анотацією помилок за типами (на основі фреймворку
MQM) та оцінкою міжфразової когерентності дає змогу виявити якісні
нюанси, недоступні для автоматичного аналізу.

\item \textbf{Розроблено класифікацію помилок перекладу} на основі
фреймворку MQM, адаптовану до специфіки дослідження. Виділено шість
основних типів помилок (пропуски, додавання, спотворення, помилки
в іменах власних, граматичні та стилістичні помилки) з трирівневою
шкалою серйозності.

\item \textbf{Описано дизайн експерименту,} що включає процедури
автоматичного та експертного оцінювання, статистичні методи аналізу
(описова статистика, кореляція Спірмена, критерій Манна-Уітні) та
заходи із забезпечення валідності дослідження.

\item \textbf{Передбачено аналіз каскадного ефекту} — впливу помилок
ASR на якість машинного перекладу через обчислення WER та його
кореляції з метриками якості перекладу.

\end{enumerate}

Представлена методологія створює основу для проведення експериментального
дослідження, результати якого наведено в наступному розділі.


\clearpage
